\section{Sharp model families and false shortcuts}
\label{sec:tpc49-models}

The weighted theorem contains two independent inputs: bounded
projective cost and admitted density.  The following exact families
show why both matter and why several qualitative substitutes fail.

\subsection{Ferrers masks spend only logarithmic energy}

Let
\begin{equation}
 T_n=(\one_{k\le r})_{1\le r,k\le n}
 \label{eq:tpc49-Ferrers-matrix}
\end{equation}
be the lower-triangular summation matrix.

\begin{proposition}[Exact Ferrers distortion]
\label{prop:tpc49-Ferrers-distortion}
The singular values of \(T_n\) are
\begin{equation}
 \sigma_k(T_n)
 =\left[2\sin\frac{(2k-1)\pi}{4n+2}\right]^{-1},
 \qquad1\le k\le n.
 \label{eq:tpc49-Ferrers-singular-values}
\end{equation}
Moreover
\begin{align}
 \|T_n\|_{S_2}^2&=\frac{n(n+1)}2,
 \label{eq:tpc49-Ferrers-Frobenius}\\
 \|T_n\|_{S_1}&=\frac n\pi\log n+O(n),
 \label{eq:tpc49-Ferrers-nuclear}\\
 r_{\rm nuc}(T_n)
 &=\frac{2}{\pi^2}(\log n)^2+O(\log n).
 \label{eq:tpc49-Ferrers-effective-rank}
\end{align}
\end{proposition}

\begin{proof}
The inverse of \(T_n\) is the lower bidiagonal first-difference matrix.
Diagonalizing \((T_n^{-1})^*T_n^{-1}\) by the discrete sine basis gives
\eqref{eq:tpc49-Ferrers-singular-values}; this is also the exact
calculation used in TPC--36.  Counting ones gives
\eqref{eq:tpc49-Ferrers-Frobenius}.  For
\(k\le n/2\), where the corresponding angle satisfies
\(0<\theta_k\le\pi/4\), use uniformly
\begin{equation}
 \frac1{2\sin\theta}
 =\frac1{2\theta}+O(\theta)
 \label{eq:tpc49-sine-expansion}
\end{equation}
and
\(\sum_{k\le n/2}(2k-1)^{-1}=\tfrac12\log n+O(1)\).
The remaining singular values sum to \(O(n)\).  This proves
\eqref{eq:tpc49-Ferrers-nuclear}; division by
\eqref{eq:tpc49-Ferrers-Frobenius} gives the last formula.
\end{proof}

TPC--36 proves \(\|T_n\|_{\pi,\infty}\asymp\log n\), while its
unweighted admitted density tends to \(1/2\).  Thus
\cref{thm:tpc49-weighted-transfer} has the correct order for this
family.  At every polynomial \(n=X^{O(1)}\), the energy cost is
\(X^{o(1)}\).  Signed Ferrers separation is not a polynomial Door-B
obstruction.

\subsection{Product kernels and product phases}

Let \(q>2\) be prime and write \(G_q=\mathbb F_q^\times\).  For
\(a\in G_q\), define
\begin{equation}
 K_a(x,y)=\one_{xy=a},
 \qquad x,y\in G_q.
 \label{eq:tpc49-product-kernel}
\end{equation}

\begin{proposition}[Cost-one product permutation]
\label{prop:tpc49-product-permutation}
The matrix \(K_a\) is a permutation of size \(q-1\), and
\begin{equation}
 \|K_a\|_{\pi,\infty}=1,
 \qquad
 r_{\rm nuc}(K_a)=q-1.
 \label{eq:tpc49-product-permutation-data}
\end{equation}
\end{proposition}

\begin{proof}
For each \(x\), the equation \(xy=a\) has the unique solution
\(y=ax^{-1}\).  Apply \cref{prop:tpc49-permutation-sharp}.  Equivalently,
multiplicative-character orthogonality gives a cost-one projective
representation.
\end{proof}

This is the same product-delta kernel that appears in the completion
analysis of TPC--36.  It is a method-specific finite-field obstruction,
not the literal residual coefficient.
At the endpoint one may choose a prime \(q\asymp J\) by Bertrand's
postulate, so the model dimension has precisely the orbit-scale exponent.

For \(r\ne0\), also define the flat product phase
\begin{equation}
 A_r(x,y)=\e_q(rxy),
 \qquad \e_q(z)=\exp(2\pi iz/q).
 \label{eq:tpc49-product-phase}
\end{equation}

\begin{proposition}[Flat entries do not imply low nuclear rank]
\label{prop:tpc49-product-phase}
The singular values of \(A_r\) are \(\sqrt q\), with multiplicity
\(q-2\), and \(1\), with multiplicity one.  Therefore
\begin{align}
 \|A_r\|_{S_2}^2&=(q-1)^2,
 \label{eq:tpc49-product-phase-Frobenius}\\
 r_{\rm nuc}(A_r)
 &=\left(
 \frac{(q-2)\sqrt q+1}{q-1}
 \right)^2\asymp q.
 \label{eq:tpc49-product-phase-rank}
\end{align}
The last displayed square is exactly the square of the bounded
projective norm computed in TPC--36.
\end{proposition}

\begin{proof}
Additive orthogonality gives
\(A_rA_r^*=qI_{q-1}-J_{q-1}\).  Its eigenvalues are \(q\) on the
orthogonal complement of constants and \(1\) on constants.  The stated
Schatten norms follow.  Multiplicative characters give flat singular
vectors and attain the bounded projective lower bound
\(\|A_r\|_{S_1}/(q-1)\); see \citep{WangTPC36}.
\end{proof}

Thus flat entries, flat singular vectors, and the usual upper-bound
notion of incoherence do not force a small projective-to-atomic loss.
They do not control the number of comparable singular values.

\subsection{Positive density alone also fails}

The permutation example has small density.  A dense nonnegative example
shows that density without bounded projective cost is equally
insufficient.  Let \(n=4^k\), put
\begin{equation}
 H_4=J_4-2I_4,
 \qquad H_n=H_4^{\otimes k},
 \qquad K_n=\frac{J_n+H_n}{2}.
 \label{eq:tpc49-dense-Hadamard-mask}
\end{equation}
Then \(H_n\) has entries \(\pm1\), satisfies
\(H_nH_n^*=nI_n\), and has constant row sum \(\sqrt n\).  Hence
\(K_n\) is a \(0\)-\(1\) matrix.

\begin{proposition}[Dense positive high-rank mask]
\label{prop:tpc49-dense-positive-mask}
The singular values of \(K_n\) are
\begin{equation}
 \frac{n+\sqrt n}{2}\quad\text{once},
 \qquad
 \frac{\sqrt n}{2}\quad\text{with multiplicity }n-1.
 \label{eq:tpc49-dense-positive-singular-values}
\end{equation}
Consequently
\begin{equation}
 \delta_w(K_n;\one,\one)
 =\frac{1+n^{-1/2}}2,
 \qquad
 r_{\rm nuc}(K_n)=\frac{n+\sqrt n}{2}.
 \label{eq:tpc49-dense-positive-data}
\end{equation}
In particular,
\(\|K_n\|_{\pi,\infty}\ge(1+\sqrt n)/2\).
\end{proposition}

\begin{proof}
Using \(J_nH_n=H_nJ_n=\sqrt nJ_n\),
\begin{equation}
 K_nK_n^*
 =\frac14\left((n+2\sqrt n)J_n+nI_n\right).
 \label{eq:tpc49-dense-positive-Gram}
\end{equation}
This gives \eqref{eq:tpc49-dense-positive-singular-values}.  The number
of ones is
\(\|K_n\|_{S_2}^2=n(n+\sqrt n)/2\), while
\(\|K_n\|_{S_1}=n(1+\sqrt n)/2\); division proves
\eqref{eq:tpc49-dense-positive-data}.  Finally
\(\|K_n\|_{\pi,\infty}\ge\|K_n\|_{S_1}/n\) by the elementary nuclear
lower bound for the bounded projective norm.
\end{proof}

The two hypotheses in \cref{thm:tpc49-weighted-transfer} are therefore
genuinely complementary: projective cost one with small density can be
bad, and positive density with polynomial projective cost can be bad.
The complete TPC static mask is favorable because it has both soft cost
and density \(1-o(1)\).
